\documentclass[11pt]{article}

\usepackage[margin=1in]{geometry}
\usepackage{amsmath, amssymb, amsthm}
\usepackage{algorithm}
\usepackage{algpseudocode}
\usepackage{hyperref}
\usepackage{enumitem}
\usepackage{booktabs}

\title{Notes on Function Approximators}
\author{}
\date{}

\begin{document}
\maketitle

\section{Motivation}

State aggregation and tile coding (see the continuous-observation-space
notes) turn a continuous state into a discrete index and keep an explicit
table entry per cell — every cell is still learned \emph{independently}.
\textbf{Function approximation} goes one step further: $\hat v$ (or
$\hat q$) is a parametric function of the state,

\[
    \hat v(s \mid \mathbf{w}) \approx v_\pi(s), \qquad \mathbf{w} \in \mathbb{R}^n,
\]

with $n$ typically far smaller than $|\mathcal{S}|$ (or, for continuous
$\mathcal{S}$, finite instead of uncountable). Updating $\mathbf{w}$ at one
state now changes the estimate at \emph{every} state whose features
overlap with it — this is exactly the generalization tile coding hinted
at, taken to its natural conclusion. The price is that $\hat v$ can no
longer represent every value function exactly; it is restricted to
whatever the parametric family can express.

\section{The prediction objective}

With a table, "best fit" was unambiguous: exact per-state values.
With function approximation, $\mathbf{w}$ has fewer dimensions than
$\mathcal{S}$, so improving the estimate at one state can worsen it at
another. We need an explicit objective. The standard choice is the
\textbf{Mean Squared Value Error}, weighted by how often each state is
visited under the data-generating policy, $\mu(s)$:

\[
    \overline{VE}(\mathbf{w}) = \sum_{s \in \mathcal{S}} \mu(s) \Big[ v_\pi(s) - \hat v(s \mid \mathbf{w}) \Big]^2.
\]

Since $\mathcal{S}$ (or $\mu$) is generally unknown or intractable to sum
over exactly, $\mathbf{w}$ is instead adjusted by \textbf{stochastic
gradient descent (SGD)} on individual, sampled states — each visited state
$S_t$ nudges $\mathbf{w}$ in the direction that reduces the squared error
\emph{at that state}, in proportion to how often such states occur under
$\mu$ automatically, simply because they are sampled with that frequency.

\section{Gradient descent for value prediction}

Suppose at time $t$ we have a target $U_t$ for $\hat v(S_t \mid
\mathbf{w})$ (the true $v_\pi(S_t)$ if it were available, or — as in the
algorithms below — a sampled/bootstrapped stand-in for it, such as the
Monte Carlo return $G_t$ or a TD target). The per-example squared error is
$\frac{1}{2}\big[U_t - \hat v(S_t \mid \mathbf{w})\big]^2$, and the SGD
update moves $\mathbf{w}$ against its gradient:

\[
    \mathbf{w}_{t+1} = \mathbf{w}_t + \alpha \Big[ U_t - \hat v(S_t \mid \mathbf{w}_t) \Big] \nabla \hat v(S_t \mid \mathbf{w}_t),
\]

where $\nabla \hat v(s \mid \mathbf{w}) = \big(\frac{\partial \hat
v(s\mid\mathbf{w})}{\partial w_1}, \dots, \frac{\partial \hat
v(s\mid\mathbf{w})}{\partial w_n}\big)^{\!\top}$. If $U_t$ itself depends
on $\mathbf{w}_t$ (as it does for TD targets, which bootstrap off
$\hat v$), this is only an approximation to the true gradient of
$\overline{VE}$ — hence \textbf{semi-gradient} methods (Section~5): they
differentiate through $\hat v(S_t \mid \mathbf{w})$ but treat the target
$U_t$ as fixed, ignoring its own dependence on $\mathbf{w}$.

\section{Linear function approximation}

The most important special case represents $\hat v$ as a \textbf{linear}
function of a hand- or auto-designed feature vector $\phi(s) \in
\mathbb{R}^n$:

\[
    f(s \mid \mathbf{w}) = \hat v(s \mid \mathbf{w}) = \mathbf{w} \cdot \phi(s)^{\top} = \sum_{i=1}^{n} w_i\, \phi_i(s).
\]

Because $\hat v$ is linear in $\mathbf{w}$, its gradient is simply the
feature vector itself, $\nabla \hat v(s \mid \mathbf{w}) = \phi(s)$, so the
general SGD update of Section~3 reduces to the especially simple

\[
    \mathbf{w}_{t+1} = \mathbf{w}_t + \alpha \Big[ U_t - \hat v(S_t \mid \mathbf{w}_t) \Big] \phi(S_t).
\]

For linear methods, $\overline{VE}$ is a quadratic (bowl-shaped) function
of $\mathbf{w}$ with a single minimum, so gradient and semi-gradient
methods have much stronger convergence guarantees than with nonlinear
approximators such as neural networks (Section~7). Everything therefore
hinges on the choice of feature map $\phi$.

\section{Constructing features: bases for linear methods}

\subsection{Polynomial basis}

Given a raw (already numeric) state vector $s = (s_1, s_2, \dots, s_k)$,
the \textbf{polynomial basis} augments the state with powers and
cross-terms of its components, letting a \emph{linear} model in the
augmented features capture nonlinear relationships in the original state:

\[
    \phi(s) = \big[\, s_1,\ s_1^2,\ \dots,\ s_1^k,\ s_2^i,\ s_3^i,\ \dots,\ s_n \,\big],
\]

with parameter vector

\[
    \mathbf{w} = [\, w_1, w_2, w_3, \dots, w_n \,],
\]

so that the approximator expands to an ordinary polynomial in the state
components:

\[
    f(s) = \hat v(s) = w_1 \cdot s_1 + w_2 \cdot s_1^2 + w_3 \cdot s_2 + \cdots
\]

Each additional power or cross-term is just one more entry of $\phi(s)$
and one more weight to learn — e.g.\ including $s_1^2$ lets the
approximator represent curvature along dimension 1 that a purely linear
term $w_1 s_1$ could not, at the cost of one extra parameter. The degree
of the polynomial (how high the powers go, and whether cross-terms like
$s_1 s_2$ are included) controls the usual expressiveness/data-efficiency
tradeoff: higher degree can fit more complex $v_\pi$, but has more weights
to estimate from the same amount of experience.

\subsection{Fourier basis}

Instead of powers, use cosines of the (rescaled-to-$[0,1]$) state:

\[
    \phi_i(s) = \cos\!\big(\pi\, \mathbf{c}_i \cdot s\big), \qquad \mathbf{c}_i \in \{0, 1, \dots, K\}^k,
\]

for integer coefficient vectors $\mathbf{c}_i$ up to some order $K$ per
dimension. Like the polynomial basis this is a fixed, hand-designed
feature map used linearly, but the Fourier basis tends to behave better
numerically (its terms are bounded in $[-1,1]$, unlike $s_1^k$ which can
blow up) and is a standard default for low-dimensional continuous control
tasks.

\subsection{Tile coding, revisited}

Tile coding (previous notes) is also linear function approximation: its
feature vector $\phi(s) \in \{0,1\}^{n \cdot \prod_i b_i}$ has exactly one
active ("$1$") entry per tiling and zero elsewhere, so
$\hat v(s\mid\mathbf{w}) = \mathbf{w}\cdot\phi(s)$ simply sums the $n$
active weights — equivalent to the "average of $n$ table lookups" view
used earlier. Unlike the polynomial or Fourier basis, tile-coding features
are \emph{local} (only nearby states share active tiles), which tends to
make learning faster and more stable, at the cost of coarser global
extrapolation.

\section{Semi-gradient TD(0) and gradient Monte Carlo}

Two concrete instantiations of the update in Section~3/4, differing only
in what stands in for $U_t$:

\begin{algorithm}[h]
\caption{Gradient Monte Carlo for $\hat v \approx v_\pi$}
\begin{algorithmic}[1]
\State Initialize $\mathbf{w}$ arbitrarily (e.g.\ $\mathbf{0}$); fix $\alpha$
\For{episode $=1,2,\dots$}
    \State Generate a full episode $S_0, A_0, R_1, \dots, S_{T-1}, A_{T-1}, R_T$ following $\pi$
    \For{$t = 0, 1, \dots, T-1$}
        \State $G_t \leftarrow$ return from step $t$ onward
        \State $\mathbf{w} \leftarrow \mathbf{w} + \alpha \big[ G_t - \hat v(S_t \mid \mathbf{w}) \big] \nabla \hat v(S_t \mid \mathbf{w})$
    \EndFor
\EndFor
\end{algorithmic}
\end{algorithm}

Since $U_t = G_t$ does not depend on $\mathbf{w}$, this is a \emph{true}
gradient method and inherits SGD's standard convergence guarantee to a
local (for linear $\hat v$, global) minimum of $\overline{VE}$.

\begin{algorithm}[h]
\caption{Semi-gradient TD(0) for $\hat v \approx v_\pi$}
\begin{algorithmic}[1]
\State Initialize $\mathbf{w}$ arbitrarily; fix $\alpha$
\For{episode $=1,2,\dots$}
    \State Initialize $S_0$
    \For{$t = 0,1,2,\dots$ until $S_t$ terminal}
        \State Take $A_t \sim \pi(\cdot \mid S_t)$, observe $R_{t+1}, S_{t+1}$
        \State $U_t \leftarrow R_{t+1} + \gamma \hat v(S_{t+1} \mid \mathbf{w})$ \quad (0 if $S_{t+1}$ terminal)
        \State $\mathbf{w} \leftarrow \mathbf{w} + \alpha \big[ U_t - \hat v(S_t \mid \mathbf{w}) \big] \nabla \hat v(S_t \mid \mathbf{w})$
    \EndFor
\EndFor
\end{algorithmic}
\end{algorithm}

Here $U_t$ bootstraps off $\hat v(S_{t+1}\mid\mathbf{w})$, so it silently
depends on the very $\mathbf{w}$ being updated — the true gradient of
$\overline{VE}$ would need to account for that dependence, but the update
above only follows $\nabla \hat v(S_t\mid\mathbf{w})$ and ignores it
(hence \emph{semi}-gradient). It converges reliably for linear $\hat v$
(to a fixed point near, though not exactly at, the $\overline{VE}$
minimum), but loses even that guarantee for general nonlinear
approximators combined with bootstrapping and off-policy data — the
well-known \textbf{deadly triad} of function approximation, bootstrapping,
and off-policy learning.

\section{Extending to control: $\hat q(s,a\mid\mathbf{w})$}

Everything above generalizes directly to action values by featurizing
state-action pairs, $\phi(s,a)$, and approximating

\[
    \hat q(s,a \mid \mathbf{w}) = \mathbf{w} \cdot \phi(s,a)^{\top},
\]

with the same semi-gradient update but plugging in a control target (e.g.
the Sarsa target $R_{t+1} + \gamma \hat q(S_{t+1}, A_{t+1} \mid
\mathbf{w})$, or the Q-learning target $R_{t+1} + \gamma
\max_{a'}\hat q(S_{t+1}, a' \mid \mathbf{w})$) in place of $U_t$, and an
$\varepsilon$-greedy (or otherwise soft) policy over $\hat q(S_t, \cdot
\mid \mathbf{w})$ in place of the tabular action-selection step. This is
the bridge from tabular/tile-coded Sarsa to modern approaches such as DQN,
where $\hat q(s,a\mid\mathbf{w})$ is a (nonlinear) neural network instead
of a linear-in-features model.

\section{Linear vs.\ nonlinear approximators}

\begin{center}
\begin{tabular}{@{}p{3.6cm}p{5.8cm}p{5.8cm}@{}}
\toprule
 & Linear (poly./Fourier/tile) & Nonlinear (neural network) \\
\midrule
Form & $\mathbf{w}\cdot\phi(s)^{\top}$, fixed hand-designed $\phi$ & $\hat v(s;\theta)$, features learned jointly with weights \\
$\overline{VE}$ landscape & convex (single minimum) & generally non-convex \\
Semi-gradient TD convergence & guaranteed (on-policy) & not guaranteed; deadly triad risk with bootstrapping + off-policy \\
Feature engineering effort & manual (choose basis, resolution) & largely automatic (representation learned) \\
Typical scale & low-dimensional continuous control & high-dimensional (pixels, large state spaces) \\
\bottomrule
\end{tabular}
\end{center}

\section{Reference}

Sutton, R.\ S., \& Barto, A.\ G.\ \emph{Reinforcement Learning: An
Introduction}, 2nd ed., Ch.~9 (On-policy Prediction with Approximation,
esp.\ 9.3--9.5 on gradient methods and polynomial/Fourier/tile-coding
bases), Ch.~10 (On-policy Control with Approximation), and
Section~11.3 (The Deadly Triad).

\end{document}
